\section{Introduction}
Counterfactual analysis~\cite{wachter2017counterfactual,mothilal2020explaining} asks a central question in data-driven decision making: \textit{``what minimal and feasible change would lead to a different outcome?''} It seeks semantically meaningful modifications that alter a prediction and provide actionable alternatives. By examining how outcomes shift under plausible changes, counterfactual reasoning supports explanation~\cite{dandl2020multi}, recourse~\cite{karimi2021algorithmic}, and auditing~\cite{lee2021auditing}, with important applications in healthcare~\cite{abrate2021counterfactual}, finance~\cite{yang2021exploring}, and law~\cite{zhang2023cfgl}. Recent work further leverages counterfactuals to improve diversity~\cite{mothilal2020explaining}, robustness~\cite{dominguez2022adversarial}, and fairness~\cite{von2022fairness}.

Most existing approaches, however, are model-centric by addressing the question ``\emph{how to change the input to fool the model}.'' Generate by perturbing features / nodes / edges
Optimizing model outputs
Tied to a specific model
Depend on gradients / decision boundary
Often produce synthetic, unrealistic instances
Limited interpretability and validity
They analyze gradients, perturb inputs, or generate synthetic edits by interrogating a predictive model~\cite{bjoru2026framework,mothilal2020explaining,dandl2020multi,pawelczyk2020learning,kanamori2020dace,lucic2021focus,vanlooveren2020interpretable,galhotra2021explaining,schleich2021geco,kavouras2023fairness}. These methods identify minimal hypothetical changes in input data, yet such edits may not correspond to real individuals or feasible interventions. 

We argue that, in high-stakes deployments~\cite{abrate2021counterfactual,yang2021exploring,lee2021auditing}, counterfactuals should instead be grounded in observed data instances. That is, the right question to ask is ``\emph{Which real alternatives explain the outcome?}'' Answer by querying existing data instances
Retrieving similar cases with different outcomes
Grounded in real data
Model-independent
Verifiable and semantically meaningful
Enable principled reasoning across modalities
Using existing examples enhances transparency and credibility, as users can directly inspect what differs and assess feasibility without relying on speculative modifications. This shifts the focus from model introspection to data organization and retrieval, framing counterfactual reasoning as a structured query over verifiable alternatives. Counterfactual analysis is indeed a database query problem.

Framing counterfactual analysis as a query problem reveals both a systems challenge and an analytic shortfall. Classical similarity search focuses on efficiently returning what is closest under a fixed metric~\cite{chaves2001searching, indyk1998approximate, malkov2020efficient}. In practice, however, users often seek the closest \emph{outcome-changing} alternative, one that clarifies decision thresholds and actionable recourse~\cite{lipton1990contrastive, miller2019explanation, shang2022whyami, bucinca2025contrastive}. This requires a contrastive query: retrieve different-outcome instances under explicit feasibility and change-cost semantics, in a form that supports interpretation. Counterfactual retrieval is therefore not a simple extension of nearest-neighbor search. It must enforce target outcomes and constraints, adopt a cost-faithful metric, and return informative rather than redundant answers~\cite{cherumanal2025diversification}, imposing concrete demands on representation, indexing, constraint checking, and set construction.

To enable effective counterfactual querying, in this paper, we introduce a two-phase framework with an \textit{offline} stage and an \textit{online} stage. In the offline stage, the system compiles admissibility rules and change-cost semantics from the data and specification space, and builds constraint-aware index structures that support efficient retrieval under these semantics. In the online stage, given a concrete query and user objectives, the system first generates feasible, low-cost candidates and then constructs a final response set according to a set-level objective. The framework explicitly separates instance-level guarantees from set-level composition. It is guided by four principles: returned alternatives must remain traceable through explicit domain-level changes; cost must faithfully reflect real transformation effort; the system must provide scalable and predictable access to admissible candidates; and the final set must satisfy collective objectives rather than merely returning individually minimal changes.


We then analyze how this framework operates across different data modalities, including tabular, temporal, graph, text, and image data. As the modality changes, admissibility checking and cost evaluation shift across semantic layers, which in turn reshapes index design, candidate generation, and validation strategies. In structured tabular data, feasibility and edits are transparent at the attribute level; in temporal, graph, or image domains, feasibility often depends on execution semantics, structural invariants, or object-level constraints, making validation a dominant system concern. We further characterize how different collective objectives shape response construction, including robustness~\cite{guyomard2023generating}, diversity~\cite{laugel2023achieving}, andfairness~\cite{kusner2017counterfactual}. These objectives elevate counterfactual 
query from simple filtering and ranking to structured subset selection, where the quality of the response depends on the joint composition of alternatives. In multi-modal settings, feasibility becomes a coordinated cross-modal validation problem, shifting the system from similarity-driven retrieval toward verification-driven processing.


% This paper argues that counterfactual querying should be standardized as a \emph{data system primitive}. We present an abstraction that organizes execution as a two-phase pipeline: (i) an offline phase that compiles admissibility and change-cost semantics and builds reusable, constraint-aware index structures over the dataset, and (ii) an online phase that, given a query, uses these structures to generate feasible, low-cost candidates and then returns a response set selected under a set-level objective. We then walk through the core principles that motivate this abstraction and specify the properties the pipeline should satisfy.

% Our goal is to provide a general framework for counterfactual querying, together with a set of principles that characterize its guarantees across tabular, time series, graph, image, and text data. Rather than prescribing a specific architecture, we formalize counterfactual querying through admissibility constraints, change-cost semantics, and collective set-level objectives, including robustness, diversity, coverage, and fairness. The framework separates instance-level guarantees from set-level composition and organizes execution into an offline–online pipeline to enable scalable and reliable processing. By distinguishing what must be guaranteed from how those guarantees are realized, we establish a principled foundation for future system designs.

% This vision paper makes five contributions. First, we formalize \emph{counterfactual querying} as a database query with explicit admissibility and change-cost semantics. Second, we propose a system framework organized as an offline–online pipeline, clarifying design choices across semantics, indexing, retrieval, and set construction. Third, we distill core principles required for interpretability, constraint faithfulness, and operational reliability. Fourth, we demonstrate how the framework generalizes across modalities, from tabular data to time series, graphs, images, and text. Finally, we characterize set-level objectives for collective counterfactual analysis, casting response construction as structured subset selection rather than simple filtering and ranking.